%% notes-tex/025-additive-baselines/sections/3-models.tex

\section{The models\secstatus{todo}}
\label{sec:models}

Let record $n$ have perturbed set $S_n$ of three genes, encoded as
$x_n \in \{0,1\}^{|G|}$ over the $|G| = 4{,}352$ genes, with label $y_n = \tau$. A
model is a set function $f : 2^G \to \mathbb{R}$, and its capacity to represent
interaction is a structural property of the hypothesis class. For distinct genes
$g$, $h$, $k$ and any $S$ disjoint from them,
%
\begin{equation}
  \Delta_{gh} f(S) = f(S \cup \{g,h\}) - f(S \cup \{g\}) - f(S \cup \{h\}) + f(S),
  \qquad
  \Delta_{ghk} f(S) = \sum_{U \subseteq \{g,h,k\}} (-1)^{3-|U|} f(S \cup U).
  \label{eq:deltas}
\end{equation}
%
A model with $\Delta_{ghk} f \equiv 0$ cannot express a three-way interaction and is
a null with respect to the label, whatever it is fit to.

\notesec{B0, train mean} $\hat{y} = \bar{y}_{\mathrm{train}}$. The zero point.

\notesec{B1, additive per-gene ridge}
$\hat{y}(S) = \beta_0 + \sum_{g \in S} \beta_g$, with $\beta$ the ridge minimizer
$(X^{\top}X + \alpha I)^{-1} X^{\top} \tilde{y}$ on the centered training labels and
$\alpha$ chosen on validation over $\{10^{-2}, \ldots, 10^{3}\}$. Both finite
differences vanish identically. B1 is the additive null of Visani, Verma and DeWitt
(bioRxiv \texttt{2026.04.23.719915}) transplanted from protein variants to gene
deletions.

\notesec{B2, additive plus recurring pairs}
B1 with one coefficient $\gamma_{gh}$ per query pair in the training split. Each
$\gamma_{gh}$ is the mean additive residual over that screen, shrunk by
$\alpha/(n_{gh} + \alpha)$ with $n_{gh}$ in the hundreds to thousands, so B2 is B1
plus one nearly unshrunk offset per screen. Second-order capacity on observed pairs,
no third-order capacity.

\notesec{B3, screen mean}
No fitted parameters. The mean training label over the records carrying the
record's query pair; if no pair in the record recurs in training, the mean of its
three gene means; failing that, the train mean. On arm R the first case always
fires. On arm Q it never does, and B3 is a mean of gene means.

\notesec{B4, query pair only}
B2 without the gene terms: a per-screen offset and nothing else. The array gene
contributes nothing, so B4's score is an upper bound on how much of the metric is
the identity of the screen. On arm Q no test record carries a pair with a
coefficient, and B4 is structurally zero.

\notesec{B5, nonlinearity on the same features}
A free 128-vector per gene, summed over the three perturbed genes and read out by a
two-hidden-layer network, trained by gradient descent with early stopping on
validation, three seeds. Exactly B1's feature space, but $\Delta_{ghk} \not\equiv 0$,
so B5 is not a null. It is the smallest departure from B1 that has interaction
capacity, and it isolates what capacity alone buys.

\notesec{The transformer}
The cell graph transformer of experiment 010, run on the 025 build with no change to
architecture, loss, optimizer or schedule (\texttt{cgt\_s0\_r\_kl\_000} for arm R,
\texttt{cgt\_s0\_q\_kl\_004} for arm Q). Its strain-dependent computation is a set
function over the three perturbed gene identities, the same input the nulls see. The
nine interaction graphs enter only through a penalty on layer-1 attention during
training, so they shape the weights and are absent from the prediction function. The
reading of the code that establishes both properties is in the 010 report and is
not repeated.

%% SOURCE: structural, derived from the model definitions above; no data involved.
\begin{table}[htbp]
\centering
\footnotesize
\caption{What each model can represent. The last two columns are properties of the
hypothesis class, not fitted results. Only B5 and the transformer can express a
three-way interaction, which is what the label is.}
\label{tab:capacity}
\begin{tabular}{llllcc}
\toprule
Model & Per-gene quantity & Set aggregation & Fit by & $\Delta_{gh} \not\equiv 0$ & $\Delta_{ghk} \not\equiv 0$ \\
\midrule
B0 & none & none & train mean & no & no \\
B1 & free scalar & sum & ridge & no & no \\
B2 & free scalar & sum, plus one offset per recurring pair & ridge & yes & no \\
B3 & label mean over training triples & mean over pairs, else genes & label means & yes & no \\
B4 & none & one offset per recurring pair & ridge & yes & no \\
B5 & free 128-vector & sum, then MLP & gradient descent & yes & yes \\
CGT & free 180-vector through an encoder & attention over $S$, mean, MLP & gradient descent & yes & yes \\
\bottomrule
\end{tabular}
\end{table}

\notesec{Comparability}
Every model is fit on the training part of the arm it is reported under, makes its
free choices on the validation part, and is reported on test. The transformer's
label normalizer is fit on the training part for arm Q and on all of S0 for arm R;
Pearson is invariant to that choice, and the nulls use the training mean only.
